\documentclass[11pt]{article}

\usepackage[margin=1in]{geometry}
\usepackage[T1]{fontenc}
\usepackage{lmodern}
\usepackage{amsmath,amssymb,amsthm}
\usepackage{mathtools}
\usepackage{enumitem}
\usepackage{microtype}
\usepackage[hidelinks]{hyperref}

\newcommand{\R}{\mathbb{R}}
\newcommand{\one}{\mathbf{1}}
\newcommand{\M}{\mathcal{M}}
\newcommand{\CC}{\mathcal{C}}
\newcommand{\UU}{\mathcal{U}}
\newcommand{\NN}{\mathcal{N}}
\newcommand{\RR}{\mathcal{R}}
\newcommand{\CD}{\mathrm{CD}}
\newcommand{\Br}{\mathrm{Br}}
\newcommand{\Acc}{\mathrm{Acc}}
\newcommand{\Del}{\mathrm{Del}}
\newcommand{\hull}{\operatorname{hull}}
\newcommand{\Rp}{\R^n_{+}}
\newcommand{\Rpp}{\inter\R^n_{+}}
\DeclareMathOperator{\conv}{conv}
\DeclareMathOperator{\cone}{cone}
\DeclareMathOperator{\relint}{relint}
\DeclareMathOperator{\rbd}{rbd}
\DeclareMathOperator{\aff}{aff}
\DeclareMathOperator{\cl}{cl}
\DeclareMathOperator{\inter}{int}

\theoremstyle{plain}
\newtheorem{theorem}{Theorem}[section]
\newtheorem{proposition}[theorem]{Proposition}
\newtheorem{lemma}[theorem]{Lemma}
\newtheorem{corollary}[theorem]{Corollary}
\theoremstyle{definition}
\newtheorem{definition}[theorem]{Definition}
\newtheorem{example}[theorem]{Example}
\theoremstyle{remark}
\newtheorem{remark}[theorem]{Remark}

\setlist{itemsep=2pt,topsep=4pt}

\title{Universal coverage by comparison deletion}
\author{Jason Konek}
\date{Draft, 26 October 2026}

\begin{document}
\maketitle

\begin{abstract}
Fix finitely many worlds and a one-sided domain $U=D_{K'}$ of gambles. A model is a set
of gambles in $U$. A ranking assigns to each world a reflexive relation on models; it is
required to be strictly guided by the truth, to compare models only where they differ,
and to satisfy a weak acyclicity condition, but not to be transitive or complete.
Starting from the Brier order, we delete, at a single world, each comparison whose record
(the pair of sets on which the two models differ) would let some model weakly dominate a
coherent model, together with the reverse comparison at that world when the two models
tie in Brier score there. Deletions are made by record, so they apply in every context;
nothing else is deleted. The resulting ranking
satisfies the requirements and finite-chain consistency, and no coherent model has a weak
dominator other than itself. Whether its undominated models are exactly the coherent ones
reduces to a geometric condition on the moment map, which we prove: every incoherent model
$X$ is improved in Brier score, by a positive margin at every world, by a coherent model
$N$ such that $N\setminus X$ lies in the smallest coherent model containing $X\setminus
N$. Consequently every incoherent model is strictly dominated by a coherent undominated
model, and the ranking covers exactly the full-dimensional credal sets $K$ with
$K'\subseteq\relint K$. These are all the credal sets that any ranking covers on $U$, so
the ranking is a universal cover on $U$; and every full-dimensional credal set is covered
on some domain. The ranking is defined by reference to coherence, need not be transitive,
and can violate disjoint additivity. A lexicographic variant satisfies disjoint
additivity and is also a universal cover on $U$, so disjoint additivity does not by
itself block universal coverage. We state these limits precisely.
\end{abstract}

\section{Introduction}\label{sec:intro}

A model is a set of gambles, the gambles it accepts. A ranking compares models at each
world; a ranking \emph{covers} a credal set $K$ if it is guided by the truth,
compares models only where they differ, makes every undominated model coherent, makes
every incoherent model dominated by a coherent undominated one, leaves $D_K$
undominated, and lives on a domain that leaves $D_K$ free to vary near its boundary
(Definition~\ref{def:cover}). The relations are required to be reflexive and to satisfy
the acyclicity condition (TC) below; transitivity and completeness are not required.

This note shows that, on a fixed one-sided domain $U=D_{K'}$, a single ranking makes the
undominated models exactly the coherent models on $U$
(Theorem~\ref{thm:main}). The ranking is obtained from the Brier order by deleting
comparisons. The deletion is arranged so that each world's relation remains a
sub-relation of the Brier order that is symmetric on Brier ties; this gives the
acyclicity condition at once (Lemma~\ref{lem:deletion}), and it makes every coherent
model immune to weak dominance (Theorem~\ref{thm:reduction}). Whether every incoherent
model is then dominated by a coherent one is equivalent to a geometric condition (H) on
the moment map and the coherent hull (Corollary~\ref{cor:H}). Theorem~\ref{thm:H} proves
(H). Section~\ref{sec:scope} records the scope of the result and its limits.

\section{Setting and framework}\label{sec:framework}

\subsection{Gambles, credal sets and the measure}

Let $\Omega=\{1,\dots,n\}$ with $n\ge2$. A \emph{gamble} is a vector $h\in\R^n$; $h_i$ is
its payoff at world $i$. Write $e_1,\dots,e_n$ for the standard basis, $\one\coloneqq
(1,\dots,1)$, $\Delta\coloneqq\{m\in\Rp:\one\cdot m=1\}$ for the simplex and
$\aff\Delta=\{m:\one\cdot m=1\}$. For $C\subseteq\aff\Delta$, $\relint C$ is the interior
of $C$ relative to $\aff\Delta$ and $\rbd C\coloneqq C\setminus\relint C$. A compact
convex $K\subseteq\Delta$ is \emph{full-dimensional} if $\relint K\neq\emptyset$. A
\emph{credal set} is a nonempty closed convex subset of $\Delta$. For $S\subseteq\R^n$,
$\cone S$ is the set of finite non-negative combinations of points of $S$, with
$\cone\emptyset=\{0\}$. We write $x\gg0$ if $x\in\Rpp$, that is, if every coordinate of
$x$ is positive.

For $S\subseteq\Delta$ define
\[
  D_S\coloneqq\{h\in\R^n: q\cdot h>0\text{ for all }q\in S\},\qquad
  C_S\coloneqq\{h\in\R^n: q\cdot h\ge0\text{ for all }q\in S\},
\]
with $D_\emptyset=C_\emptyset\coloneqq\R^n$, and, for nonempty compact $S$,
$P_S(h)\coloneqq\min_{q\in S}q\cdot h$.

\begin{itemize}
  \item[(M)] $\mu$ is a finite Borel measure on $\R^n$ with the same null sets as Lebesgue
  measure, and $\int|h|\,d\mu(h)<\infty$.
\end{itemize}
Throughout, $\mu$ satisfies (M). Equalities and inclusions between sets of gambles hold
\emph{almost everywhere} (a.e.)
if they hold up to a $\mu$-null set. Every hyperplane is null, and a nonempty open set is
never null. For measurable $X$ put
\[
  m(X)\coloneqq\int_Xh\,d\mu(h)\in\R^n,\qquad m_i(X)\coloneqq\int_Xh_i\,d\mu(h).
\]
For measurable $E$, $E_d$ denotes the set of Lebesgue density points of $E$. By the
Lebesgue density theorem $E\setminus E_d$ is null, and by (M), $\mu(E\cap V)>0$ for
every open ball $V$ centred at a point of $E_d$.

\begin{lemma}\label{lem:D}
Let $S,S_1,S_2\subseteq\Delta$ be nonempty and compact.
\begin{enumerate}[label=(\alph*)]
  \item $D_S$ is an open convex cone containing $\Rpp$, and $D_S\cap(-D_S)=\emptyset$.
  Moreover $D_S=D_{\conv S}$ and $D_{S_1}\cap D_{S_2}=D_{S_1\cup S_2}$.
  \item $P_S(h+t\one)=P_S(h)+t$ for $t\in\R$. The boundary $\partial D_S$ is contained in
  $\{P_S=0\}$, and $\{P_S=0\}$ is null.
  \item If $S_1,S_2$ are convex and $D_{S_1}\subseteq D_{S_2}$, then $S_2\subseteq S_1$.
  \item If $O$ is open and $O\subseteq D_S$ a.e., then $O\subseteq D_S$.
\end{enumerate}
\end{lemma}

\begin{proof}
(a) $|P_S(h)-P_S(h')|\le\max_j|h_j-h'_j|$ because $S\subseteq\Delta$, so $D_S=\{P_S>0\}$
is open. Convexity, the cone property and $\Rpp\subseteq D_S$ are immediate, and $q\cdot
h>0$, $q\cdot(-h)>0$ cannot both hold. A linear function has the same minimum over $S$
and over $\conv S$. The last identity is the definition.
(b) The identity holds because $q\cdot\one=1$ for $q\in\Delta$. If $P_S(h)<0$ then $h$
has a neighbourhood outside $D_S$; if $P_S(h)=0$ then $h+t\one\in D_S$ for $t>0$ and
$h\notin D_S$. So $\partial D_S\subseteq\{P_S=0\}$. Write $h=y+t\one$ with $y\in
\one^\perp$, a linear change of coordinates. Then $\{P_S=0\}=\{y-P_S(y)\one:y\in
\one^\perp\}$ meets each line $y+\R\one$ in one point, so it is Lebesgue-null by Fubini's
theorem, and $\mu$-null by (M).
(c) Suppose $k\in S_2\setminus S_1$. Strict separation of $k$ from $S_1$ gives $g\in\R^n$
and $c\in\R$ with $g\cdot q>c>g\cdot k$ for all $q\in S_1$. Put $h\coloneqq g-c\one$. Then
$q\cdot h>0$ on $S_1$ and $k\cdot h<0$, so $h\in D_{S_1}\setminus D_{S_2}$.
(d) The open set $O\cap\{P_S<0\}$ is null, hence empty. If $h\in O$, then $h-\varepsilon
\one\in O$ for some $\varepsilon>0$, so $P_S(h)-\varepsilon\ge0$ by (b), and $h\in D_S$.
\end{proof}

\subsection{Models, coherence and the coherent hull}

A \emph{domain} is a measurable set $U\subseteq\R^n$. The \emph{models} on $U$ are the
measurable sets $X\subseteq U$, identified when they agree a.e.; their class is $\M_U$.
A model $X$ \emph{accepts} $h$ if $h\in X$. At world $i$, accepting $h$ is \emph{correct}
if $h_i>0$ and \emph{wrong} if $h_i<0$. A model $X$ is \emph{coherent} if $X=D_Q$ a.e.\
for some credal set $Q$. If $D_{Q_1}=D_{Q_2}$ a.e., Lemma~\ref{lem:D}(d) gives $D_{Q_1}=
D_{Q_2}$, and Lemma~\ref{lem:D}(c) gives $Q_1=Q_2$; so $Q$ is unique. We call $U$
\emph{one-sided} if $U\cap(-U)$ is null. This usage differs from that of
\cite{KCM}, where ``one-sided'' means $U\cap(-U)=\emptyset$ and ``essentially one-sided''
means that $U\cap(-U)$ contains no non-null set of gambles with a positive and a negative
coordinate; our condition lies between the two.

\smallskip\noindent\textbf{Standing assumption.} Except where a general domain is
mentioned explicitly, $K'\subseteq\Delta$ is nonempty, compact and convex, and
$U\coloneqq D_{K'}$. By Lemma~\ref{lem:D}(a), $U$ is open, $U\supseteq\Rpp$ and $U\cap
(-U)=\emptyset$; moreover $U$ contains no $h$ with $h\le0$, since $k\cdot h\le0$ for
$k\in K'$.

\smallskip
Let $\CC_U$ be the class of coherent models on $U$. If $Q$ is a credal set with
$K'\subseteq Q$, then $D_Q\subseteq D_{K'}=U$, so $D_Q\in\CC_U$. Conversely, if $D_Q
\subseteq U$ a.e., then $D_Q\subseteq D_{K'}$ by Lemma~\ref{lem:D}(d), and $K'\subseteq
Q$ by Lemma~\ref{lem:D}(c). Hence
\[
  \CC_U=\{D_Q: Q\text{ a credal set with }K'\subseteq Q\}.
\]
Every coherent model contains $\Rpp$. For measurable $B\subseteq U$ put
\[
  Q_B\coloneqq\{q\in\Delta: q\cdot h\ge0\text{ for $\mu$-a.e.\ }h\in B\},\qquad
  \hull(B)\coloneqq D_{Q_B}.
\]

\begin{lemma}\label{lem:hull}
Let $B\subseteq U$ be measurable.
\begin{enumerate}[label=(\alph*)]
  \item $Q_B$ is a credal set and $K'\subseteq Q_B$. Hence $\hull(B)\in\CC_U$.
  \item $B\subseteq\hull(B)$ a.e.
  \item For a credal set $Q$ with $K'\subseteq Q$, the following are equivalent:
  $B\subseteq D_Q$ a.e.; $Q\subseteq Q_B$; $\hull(B)\subseteq D_Q$. So $\hull(B)$ is the
  smallest coherent model containing $B$ a.e.
\end{enumerate}
\end{lemma}

\begin{proof}
(a) $K'\subseteq Q_B$ because $k\cdot h>0$ on $U$ for $k\in K'$. Convexity is clear. Let
$q_k\in Q_B$ converge to $q$. If $q\cdot h<0$ then $q_k\cdot h<0$ for all large $k$, so
$B\cap\{q\cdot h<0\}\subseteq\bigcup_kB\cap\{q_k\cdot h<0\}$, which is null. So $Q_B$ is
closed. Since $Q_B\supseteq K'$, $\hull(B)\in\CC_U$.
(b) $B\setminus D_{Q_B}\subseteq\{P_{Q_B}=0\}\cup\bigcup_{q\in Q_B}\{q\cdot h<0\}$. The
first set is null by Lemma~\ref{lem:D}(b). The open sets $\{q\cdot h<0\}$ have a countable
subcover $\{q_k\cdot h<0\}$ (Lindel\"of), and each $B\cap\{q_k\cdot h<0\}$ is null.
(c) If $B\subseteq D_Q$ a.e., every $q\in Q$ has $q\cdot h>0$ for a.e.\ $h\in B$, so
$Q\subseteq Q_B$. If $Q\subseteq Q_B$, then $D_{Q_B}\subseteq D_Q$ by definition. If
$\hull(B)\subseteq D_Q$, then $B\subseteq D_Q$ a.e.\ by (b).
\end{proof}

\subsection{Rankings and dominance}

A \emph{ranking} on $\M_U$ is a family $(\preceq_i)_{i\in\Omega}$ of binary relations on
$\M_U$. Lower is better: $X\preceq_iY$ says that $X$ does at least as well as $Y$ at
world $i$. We write $X\prec_iY$ if $X\preceq_iY$ and not $Y\preceq_iX$; $X$ and $Y$ are
\emph{tied} at $i$ if $X\preceq_iY$ and $Y\preceq_iX$, and \emph{incomparable} at $i$ if
neither holds. We assume throughout that each $\preceq_i$ is reflexive, that a.e.-equal
models are tied, and
\begin{itemize}
  \item[(TC)] if $A\preceq_iB$ and $B\preceq_iC$, then not $C\prec_iA$.
\end{itemize}
Transitivity and completeness are not assumed. We also consider the stronger condition
\begin{itemize}
  \item[(FCC)] \emph{Finite-chain consistency.} If $A_0\preceq_iA_1\preceq_i\dots
  \preceq_iA_k$, then not $A_k\prec_iA_0$.
\end{itemize}
(TC) is the case $k=2$ of (FCC).

\begin{definition}\label{def:dom}
Let $X,Y\in\M_U$. $Y$ \emph{weakly dominates} $X$ if $Y\preceq_iX$ for every $i$. $Y$
\emph{dominates} $X$ if $Y$ weakly dominates $X$ and $Y\prec_iX$ for some $i$. $Y$
\emph{strictly dominates} $X$ if $Y\prec_iX$ for every $i$. $X$ is \emph{undominated} if
no model dominates it. We write $\UU$ for the undominated models, $\CC$ for the coherent
models and $\CC^\ast\coloneqq\CC\cap\UU$.
\end{definition}

If $Y$ weakly dominates $X$ but does not dominate it, then $X$ and $Y$ are tied at every
world; this uses only the definition of $\prec_i$.

The \emph{Brier loss} of $X\in\M_U$ at world $i$ is the $|h_i|$-weighted measure of its
wrong acceptances and of its silences on correct gambles:
\[
  \beta_i(X)\coloneqq\int_{X\cap\{h_i<0\}}|h_i|\,d\mu
  +\int_{(U\setminus X)\cap\{h_i>0\}}|h_i|\,d\mu=c_i-m_i(X),\qquad
  c_i\coloneqq\int_{U\cap\{h_i>0\}}h_i\,d\mu .
\]
The \emph{Brier ranking} is $X\preceq^{\mathrm B}_iY$ iff $\beta_i(X)\le\beta_i(Y)$, that
is, iff $m_i(X)\ge m_i(Y)$. $Y$ \emph{uniformly dominates} $X$ in Brier loss if there is
$\varepsilon>0$ with $\beta_i(Y)+\varepsilon\le\beta_i(X)$ for every $i$, that is, if
$m(Y)-m(X)\gg0$.

\subsection{Requirements}

\begin{itemize}
  \item[(G)] \emph{Strict guidance.} If $X\setminus Y\subseteq\{h_i>0\}$ a.e.,
  $Y\setminus X\subseteq\{h_i<0\}$ a.e., and $X\triangle Y$ is not null, then
  $X\prec_iY$.
  \item[(IB)] \emph{Independence.} Let $S$ be measurable and $X,Y,X',Y'\in\M_U$ with
  $X\setminus S=Y\setminus S$, $X'\setminus S=Y'\setminus S$, $X\cap S=X'\cap S$ and
  $Y\cap S=Y'\cap S$, all a.e. Then $X\preceq_iY$ if and only if $X'\preceq_iY'$.
  \item[(GC)] \emph{Global coherence.} Every undominated model is coherent.
  \item[(C)] \emph{Direct closure under ties.} A model that is tied at every world with a
  member of $\CC^\ast$ belongs to $\CC^\ast$.
  \item[(ID)] \emph{Incoherent-model dominance.} Every incoherent model is dominated by a
  member of $\CC^\ast$.
  \item[(LF)] \emph{Local freedom at $D_K$.} There is an open set $N\supseteq\partial
  D_K\setminus\{0\}$ such that every measurable $X$ with $X\triangle D_K\subseteq N$ a.e.\
  belongs to $\M_U$, that is, satisfies $X\subseteq U$ a.e.
  \item[(W)] Every model is weakly dominated by a member of $\CC^\ast$.
\end{itemize}

\begin{definition}\label{def:cover}
Let $K$ be a credal set and $U$ a domain. A ranking on $\M_U$ \emph{covers} $K$ if:
\begin{enumerate}[label=(\roman*)]
  \item it satisfies (G) and (IB);
  \item it satisfies (GC);
  \item it satisfies (C);
  \item $D_K\in\M_U$ and $D_K$ is undominated;
  \item it satisfies (ID);
  \item (LF) holds at $D_K$.
\end{enumerate}
\end{definition}

(W) is not part of the definition; the ranking constructed below satisfies it.

\begin{lemma}\label{lem:LF}
Let $K$ be a credal set.
\begin{enumerate}[label=(\alph*)]
  \item $D_K\in\M_U$ if and only if $K'\subseteq K$.
  \item If $K$ is full-dimensional and $K'\subseteq K$, then (LF) holds at $D_K$ if and
  only if $K'\subseteq\relint K$.
  \item Let $V$ be any one-sided domain. If $K$ is not full-dimensional, then (LF) fails
  at $D_K$ for $\M_V$.
\end{enumerate}
\end{lemma}

\begin{proof}
(a) This is the description of $\CC_U$ above.

(b) First let $h\in C_K\setminus\{0\}$ and $x\in\relint K$. If $x\cdot h=0$, then for
every $d$ with $\one\cdot d=0$ and small $t>0$ both $x\pm td$ lie in $K$, so $\pm t\,d
\cdot h\ge0$ and $d\cdot h=0$. Then $h=c\one$ for some $c$, and $c=x\cdot h=0$, a
contradiction. So $x\cdot h>0$.

Suppose $K'\subseteq\relint K$. Then $D_K\subseteq U$. Each $h\in\partial D_K\setminus
\{0\}$ lies in $C_K\setminus\{0\}$, so $k\cdot h>0$ for $k\in K'$ by the previous
paragraph, that is, $h\in U$. Take $N\coloneqq U$, which is open. If $X\triangle D_K
\subseteq U$ a.e., then $X\subseteq D_K\cup U=U$ a.e.

Conversely, suppose (LF) holds with the open set $N$, and $x\in K'\cap\rbd K$. Since
$x\notin\relint K$, which is the relative interior of $K$ in the sense of convex analysis
because $K$ is full-dimensional, there is a nontrivial supporting hyperplane to $K$
through $x$ \cite[Thm.~11.6]{Rockafellar}: some $g\in\R^n$ has $g\cdot k\ge g\cdot x$ for
all $k\in K$, with strict inequality for some $k\in K$. Put $h\coloneqq g-(g\cdot x)\one$.
Then $k\cdot h\ge0$ on $K$, $x\cdot h=0$ and $h\ne0$. So $h\in C_K\setminus D_K$, and
$h+t\one\in D_K$ for $t>0$; hence $h\in\partial D_K\setminus\{0\}\subseteq N$. The model
$D_K\cup N$ satisfies $(D_K\cup N)\triangle D_K\subseteq N$, so it belongs to $\M_U$;
thus $N\subseteq U$ a.e., and $N\subseteq U$ by Lemma~\ref{lem:D}(d). But $x\cdot h=0$
with $x\in K'$, so $h\notin U$, a contradiction.

(c) The linear span of $K$ has dimension at most $n-1$, so some $h\ne0$ has $k\cdot h=0$
for all $k\in K$. Then $\pm h\in C_K\setminus D_K$ and $\pm h+t\one\in D_K$ for $t>0$, so
$\pm h\in\partial D_K\setminus\{0\}$. If (LF) held with the open set $N$, then, as in (b),
$N\subseteq V$ a.e.; choose $r>0$ with $B(\pm h,r)\subseteq N$. Then $B(h,r)\subseteq V$
a.e.\ and $B(h,r)=-B(-h,r)\subseteq-V$ a.e., so the non-null ball $B(h,r)$ is contained
in $V\cap(-V)$ a.e., which contradicts one-sidedness.
\end{proof}

\section{The deletion construction}\label{sec:deletion}

\subsection{Records}

A \emph{record} is an ordered pair $r=(A,B)$ of disjoint measurable subsets of $U$, with
records identified when their components agree a.e. The record of a pair of models is
\[
  r(X,Y)\coloneqq(X\setminus Y,\;Y\setminus X).
\]
It reads: $X$ accepts $A$ where $Y$ accepts $B$. The \emph{inverse} of $r=(A,B)$ is
$r^{-1}\coloneqq(B,A)$, so $r(Y,X)=r(X,Y)^{-1}$. The record $r$ is \emph{null} if $A$
and $B$ are null, and its \emph{moment} is $m(r)\coloneqq m(A)-m(B)$; thus $m(r(X,Y))=
m(X)-m(Y)$. A record $r=(A,B)$ \emph{attacks} a model $Z$ if $B\subseteq Z$ a.e.\ and
$A\cap Z$ is null. Then $Z\oplus r\coloneqq(Z\setminus B)\cup A$ satisfies $r(Z\oplus r,
Z)=r$: the record $r(N,X)$ attacks $Z$ exactly when the change that turns $X$ into $N$
can be applied to $Z$. In particular $r(Y,Z)$ attacks $Z$ for all $Y,Z\in\M_U$.

\begin{lemma}\label{lem:record}
Let $\Acc$ be a set of records containing the null record, and define $X\preceq Y$ iff
$r(X,Y)\in\Acc$. Then $\preceq$ is reflexive, ties a.e.-equal models and satisfies (IB).
\end{lemma}

\begin{proof}
If $X=Y$ a.e., then $r(X,Y)$ and $r(Y,X)$ are null. For (IB), the hypotheses give
$X\triangle Y\subseteq S$ a.e., so $X\setminus Y=(X\cap S)\setminus(Y\cap S)=(X'\cap S)
\setminus(Y'\cap S)=X'\setminus Y'$ a.e., and likewise $Y\setminus X=Y'\setminus X'$
a.e. So $r(X,Y)=r(X',Y')$.
\end{proof}

\subsection{The deletion lemma}

\begin{lemma}\label{lem:deletion}
Let $\preceq$ be a relation on a set $\M$, let $(T,\le)$ be a totally ordered set, and
let $\varphi:\M\to T$ satisfy
\begin{enumerate}[label=(\alph*)]
  \item if $X\preceq Y$ then $\varphi(X)\le\varphi(Y)$;
  \item if $X\preceq Y$ and $\varphi(X)=\varphi(Y)$, then $Y\preceq X$.
\end{enumerate}
Then $\preceq$ satisfies (FCC), and in particular (TC).
\end{lemma}

\begin{proof}
If $X\prec Y$, then $\varphi(X)<\varphi(Y)$: otherwise $\varphi(X)=\varphi(Y)$ by (a), and
(b) gives $Y\preceq X$. Let $A_0\preceq A_1\preceq\dots\preceq A_k$. By (a),
$\varphi(A_0)\le\varphi(A_k)$, so $A_k\prec A_0$, which would need $\varphi(A_k)<
\varphi(A_0)$, is impossible.
\end{proof}

In words: a relation contained in the weak order of a score with values in a totally
ordered set, and symmetric on the
ties of that score, satisfies (FCC). Deleting comparisons from a score order, while
keeping the deletions symmetric on score ties, therefore never breaks (TC).

\subsection{Comparison-deleted Brier rankings}

Fix a family $\NN\subseteq\CC_U$ of coherent models. Put
\begin{align*}
  \Br_i&\coloneqq\{r: m_i(r)\ge0\},\\
  G_i&\coloneqq\{r=(A,B): A\subseteq\{h_i>0\}\text{ and }B\subseteq\{h_i<0\}\text{
  a.e.}\},\\
  \RR(\NN)&\coloneqq\{r\text{ non-null}: m(r)\in\Rp\text{ and $r$ attacks some member of
  }\NN\}.
\end{align*}
$\Br_i$ is the set of records accepted by the Brier ranking at world $i$, and $G_i$ is
the set of records that (G) concerns. For $r\in\RR(\NN)$, the \emph{deletion world}
$i(r)$ is the least $i$ with $m_i(r)>0$ and $r\notin G_i$ if there is such an $i$, and
otherwise the least $i$ with $r\notin G_i$.

The deletion world exists. Suppose $r=(A,B)\in\RR(\NN)$ lies in $G_i$ for every $i$. Then
$A\subseteq\Rpp$ a.e.\ and $B\subseteq\{h:h_i<0\text{ for all }i\}$ a.e. Since $U$
contains no such $h$,
$B$ is null. Since $r$ attacks some $D_Q\in\NN$ and $D_Q\supseteq\Rpp$, $A=A\cap D_Q$
a.e.\ is null. So $r$ is null, a contradiction.

\begin{definition}\label{def:CD}
The \emph{comparison-deleted Brier ranking} $\CD(\NN)$ on $\M_U$ is given by
\begin{align*}
  \Del_i&\coloneqq\{r\in\RR(\NN): i(r)=i\}\cup\{r^{-1}: r\in\RR(\NN),\ i(r)=i,\
  m_i(r)=0\},\\
  \Acc_i&\coloneqq\Br_i\setminus\Del_i,\qquad
  X\preceq_iY\iff r(X,Y)\in\Acc_i .
\end{align*}
\end{definition}

Each deletion depends only on the record, so it applies in every context. For
$N\in\NN$ and a model $Y\ne N$ with $m(Y)\ge m(N)$ coordinatewise, the record
$r(Y,N)$ attacks $N$, so it belongs to $\RR(\NN)$, and the comparison $Y\preceq_iN$ is
removed at the world $i=i(r(Y,N))$; if $m_i(Y)=m_i(N)$ there, the reverse comparison is
removed too. At that world $Y$ and $N$ become incomparable. No other comparison is
removed, and records in $G_i$ are never removed at world $i$
(Theorem~\ref{thm:reduction}(a)).

\section{The reduction theorem}\label{sec:reduction}

For $X\in\M_U$ and $N\in\NN$ consider the condition
\begin{itemize}
  \item[($\star$)] $m(N)\ge m(X)$ coordinatewise, $m(N)\ne m(X)$, and $r(N,X)$ attacks
  no member of $\NN$.
\end{itemize}

\begin{theorem}\label{thm:reduction}
Let $\NN\subseteq\CC_U$, and consider $\CD(\NN)$.
\begin{enumerate}[label=(\alph*)]
  \item Each $\preceq_i$ is reflexive, ties a.e.-equal models, and satisfies (G), (IB),
  (FCC) and (TC). Moreover $X\preceq_iY$ implies $m_i(X)\ge m_i(Y)$, and if in addition
  $m_i(X)=m_i(Y)$ then $Y\preceq_iX$.
  \item No member of $\NN$ has a weak dominator other than itself. In particular every
  member of $\NN$ is undominated, and a model tied at every world with a member of $\NN$
  equals it.
  \item If $X\in\M_U$ and $N\in\NN$ satisfy ($\star$), then $N$ dominates $X$. If moreover
  $m(N)-m(X)\gg0$, then $N$ strictly dominates $X$.
  \item If $N\in\NN$ dominates $X\in\M_U$, then $X$ and $N$ satisfy ($\star$).
  \item The following are equivalent: (1) every $X\in\M_U\setminus\NN$ satisfies
  ($\star$) with some $N\in\NN$; (2) every model outside $\NN$ is dominated by a member of
  $\NN$. If they hold, then $\UU=\CC^\ast=\NN$, and $\CD(\NN)$ satisfies (GC), (C) in the
  strong form of (b), (ID) and (W).
\end{enumerate}
\end{theorem}

\begin{proof}
(a) The null record lies in $\Br_i$. It is not in $\RR(\NN)$, whose members are
non-null, and it is not the inverse of a member of $\RR(\NN)$. So it lies in $\Acc_i$,
and Lemma~\ref{lem:record} gives reflexivity, ties of a.e.-equal models and (IB).

\emph{(G).} Let $r=(A,B)\in G_i$ be non-null. Since $\{h_i=0\}$ is null,
$m_i(r)=\int_A|h_i|\,d\mu+\int_B|h_i|\,d\mu>0$, so $r\in\Br_i$ and $r^{-1}\notin\Br_i$.
The record $r$ is not deleted at $i$: it is not a member of $\RR(\NN)$ with deletion
world $i$, because $r\notin G_{i(r)}$ for every $r\in\RR(\NN)$; and it is not an inverse
$s^{-1}$ with $m_i(s)=0$, because $m_i(r)\ne0$. So $r\in\Acc_i$ and $r^{-1}\notin
\Acc_i$; that is, $X\prec_iY$ whenever $r(X,Y)=r$.

\emph{Brier compatibility.} $\Acc_i\subseteq\Br_i$, and $m_i(r(X,Y))=m_i(X)-m_i(Y)$.

\emph{Tie symmetry.} Let $r\in\Acc_i$ with $m_i(r)=0$, and suppose $r^{-1}\notin
\Acc_i$. Since $m_i(r^{-1})=0$, $r^{-1}\in\Br_i$, so $r^{-1}\in\Del_i$. If $r^{-1}\in
\RR(\NN)$ with $i(r^{-1})=i$, then, as $m_i(r^{-1})=0$, its inverse $r$ lies in $\Del_i$.
If $r^{-1}=s^{-1}$ with $s\in\RR(\NN)$ and $i(s)=i$, then $s=r$ and $r\in\Del_i$. Both
contradict $r\in\Acc_i$.

\emph{(FCC) and (TC).} Apply Lemma~\ref{lem:deletion} with $\varphi\coloneqq-m_i$, that
is, with the Brier loss $\beta_i$ up to a constant.

(b) Let $Y$ weakly dominate $N\in\NN$ with $Y\ne N$ a.e., and put $r\coloneqq r(Y,N)$.
Then $r\in\Acc_i\subseteq\Br_i$ for every $i$, so $m(r)\in\Rp$; $r$ is non-null; and $r$
attacks $N$. Hence $r\in\RR(\NN)$ and $r\in\Del_{i(r)}$, contradicting $r\in
\Acc_{i(r)}$. A model tied at every world with $N$ weakly dominates $N$, so it equals $N$
a.e.

(c) Let $d\coloneqq r(N,X)$. Then $d\in\Br_i$ for every $i$. It is not in $\RR(\NN)$,
since it attacks no member of $\NN$. It is not deleted as an inverse: that would need
$d^{-1}=r(X,N)\in\RR(\NN)$, hence $m(X)\ge m(N)$, hence $m(X)=m(N)$, which is excluded. So
$d\in\Acc_i$ and $N\preceq_iX$ for every $i$. For a world $j$ with $m_j(N)>m_j(X)$,
$r(X,N)\notin\Br_j$, so not $X\preceq_jN$, and $N\prec_jX$. If $m(N)-m(X)\gg0$, this
holds at every $j$.

(d) Let $d\coloneqq r(N,X)$. Since $N\preceq_iX$ for every $i$, $d\in\Br_i$ for every
$i$, so $m(N)\ge m(X)$. Since $N$ dominates $X$, $N\ne X$ a.e., so $d$ is non-null. If $d$
attacked a member of $\NN$, then $d\in\RR(\NN)$ and $d\in\Del_{i(d)}$, contradicting
$d\in\Acc_{i(d)}$. If $m(N)=m(X)$, then $m_i(d)=0$ for every $i$, and tie symmetry gives
$X\preceq_iN$ for every $i$, so $N$ would not dominate $X$.

(e) (1)$\Rightarrow$(2) is (c), and (2)$\Rightarrow$(1) is (d). Assume (2). By (b),
$\NN\subseteq\UU$, and by (2), $\UU\subseteq\NN$; since $\NN\subseteq\CC$,
$\UU=\CC^\ast=\NN$. (GC) follows. (C) holds because a model tied at every world with a
member of $\NN$ equals it, by (b). An incoherent model is not in $\NN$, so (2) gives (ID).
For (W): a member of $\NN$ weakly dominates itself, and every other model is dominated by
a member of $\NN$.
\end{proof}

We now take $\NN=\CC_U$, all coherent models on $U$, and translate ($\star$) into a
condition on the coherent hull.

\begin{itemize}
  \item[(H)] Every incoherent $X\in\M_U$ has a coherent $N\in\M_U$ with $m(N)\ge m(X)$
  coordinatewise, $m(N)\ne m(X)$, and $\mu\bigl((N\setminus X)\cap\hull(X\setminus N)
  \bigr)>0$.
\end{itemize}

\begin{corollary}\label{cor:H}
Let $X\in\M_U$ and $N\in\CC_U$. The record $r(N,X)$ attacks some coherent model if and only
if $(N\setminus X)\cap\hull(X\setminus N)$ is null. Consequently $\CD(\CC_U)$ satisfies
(ID) if and only if (H) holds; and if (H) holds, all conclusions of
Theorem~\ref{thm:reduction}(e) hold for $\NN=\CC_U$.
\end{corollary}

\begin{proof}
Write $A\coloneqq N\setminus X$ and $B\coloneqq X\setminus N$. If $A\cap\hull(B)$ is
null, then $r(N,X)=(A,B)$ attacks $\hull(B)$, which is a coherent model containing $B$
a.e.\ (Lemma~\ref{lem:hull}(a),(b)). Conversely, if $(A,B)$ attacks a coherent model
$D_Q$, then $B\subseteq D_Q$ a.e., so $\hull(B)\subseteq D_Q$ by
Lemma~\ref{lem:hull}(c), and $A\cap\hull(B)\subseteq A\cap D_Q$ is null.

So, for $\NN=\CC_U$, condition ($\star$) for $X$ and $N$ is exactly the condition on $N$
in (H). If (ID) holds, every incoherent $X$ is dominated by a member of $\CC^\ast\subseteq
\CC_U$, and Theorem~\ref{thm:reduction}(d) gives ($\star$), hence (H) at $X$. If (H)
holds, condition (1) of Theorem~\ref{thm:reduction}(e) holds, because the models outside
$\CC_U$ are the incoherent ones.
\end{proof}

The hull condition in (H) has a finite geometric form.

\begin{lemma}\label{lem:hullformula}
For measurable $B\subseteq U$,
\[
  \hull(B)=\cone(B_d)+\Rpp .
\]
\end{lemma}

\begin{proof}
First, if $q\in Q_B$ and $b\in B_d$, then $q\cdot b\ge0$: otherwise an open ball $V$
around $b$ lies in $\{q\cdot h<0\}$, and $\mu(B\cap V)>0$ contradicts $q\in Q_B$.
Conversely, if $q\in\Delta$ and $q\cdot b\ge0$ for all $b\in B_d$, then $q\in Q_B$,
because a.e.\ point of $B$ lies in $B_d$.

Put $E\coloneqq\cone(B_d)+\Rpp$, a nonempty open convex set. If $a=\sum_j\lambda_jb_j+p$
with $b_j\in B_d$, $\lambda_j\ge0$ and $p\gg0$, then every $q\in Q_B$ has $q\cdot a\ge
q\cdot p>0$; so $a\in D_{Q_B}$. Now let $a\notin E$. Separating $a$ from $E$ gives
$q\ne0$ with $q\cdot a\le q\cdot e$ for all $e\in E$. Then $q\ge0$, since otherwise
$e=\one+se_i$ with $q_i<0$ makes $q\cdot e\to-\infty$ as $s\to\infty$. For $b\in B_d$
and $s,t>0$, $sb+t\one\in E$; letting $t\to0$ gives $q\cdot a\le s\,q\cdot b$ for all
$s>0$, so $q\cdot b\ge0$ (let $s\to\infty$); and $t\one\in E$ gives $q\cdot a\le0$. Then $\hat q\coloneqq
q/(\one\cdot q)\in\Delta$ is non-negative on $B_d$, so $\hat q\in Q_B$, and $\hat q\cdot
a\le0$; hence $a\notin D_{Q_B}$.
\end{proof}

\begin{corollary}\label{cor:witness}
Let $A,B\subseteq U$ be measurable. Then $\mu(A\cap\hull(B))>0$ if and only if there are
$a\in A_d$, points $b_1,\dots,b_r\in B_d$ with $0\le r\le n$, and $\lambda_j\ge0$, such
that $a-\sum_{j=1}^r\lambda_jb_j\gg0$.
\end{corollary}

\begin{proof}
$\hull(B)$ is open. If $a\in A_d$ lies in $\hull(B)$, a ball around $a$ lies in
$\hull(B)$ and meets $A$ in positive measure. Conversely, if $\mu(A\cap\hull(B))>0$,
some $a\in A_d$ lies in $\hull(B)$, since $A\setminus A_d$ is null. By
Lemma~\ref{lem:hullformula}, $a=\sum_j\lambda_jb_j+p$ with $b_j\in B_d$, $\lambda_j\ge0$
and $p\gg0$, and by the conic form of Carath\'eodory's theorem the sum can be rewritten
with at most $n$ linearly independent $b_j$.
\end{proof}

So (H) asks for a coherent Brier improvement $N$ of $X$ that adds some gamble $a$ which
coherence already requires of finitely many gambles that $N$ removes from $X$.

\section{The geometric condition}\label{sec:H}

For a model $X\in\M_U$ write
\[
  G\coloneqq\hull(X)=D_{Q_X},\qquad P\coloneqq P_{Q_X},\qquad
  Q_\tau\coloneqq(1-\tau)Q_X+\tau\Delta\quad(0<\tau<1).
\]
$Q_\tau$ is a compact convex subset of $\Delta$, and
\begin{equation}\label{eq:Ptau}
  P_{Q_\tau}(h)=(1-\tau)P(h)+\tau\min_ih_i .
\end{equation}
By Lemma~\ref{lem:hull}(b), $X\subseteq G$ a.e., so $P>0$ a.e.\ on $X$.

\begin{lemma}\label{lem:contraction}
Let $X\in\M_U$, $c\ge0$ and $q\in\Delta$. If $q\cdot x\ge-c\,P(x)$ for $\mu$-a.e.\ $x\in
X$, then $q\in Q_X$.
\end{lemma}

\begin{proof}
For $q'\in Q_X$ put $T(q')\coloneqq(q+cq')/(1+c)\in\Delta$. For a.e.\ $x\in X$,
$(q+cq')\cdot x\ge q\cdot x+c\,P(x)\ge0$, so $T(q')\in Q_X$. Thus $T$ maps $Q_X$ into
itself, and it is a contraction with ratio $c/(1+c)<1$ and fixed point $q$. For any
$q'\in Q_X$ (which is nonempty), $T^k(q')\to q$, and $Q_X$ is closed
(Lemma~\ref{lem:hull}(a)); so $q\in Q_X$.
\end{proof}

\begin{corollary}\label{cor:shrink}
Let $X\in\M_U$ and $0<\tau<1$. For every measurable $N\subseteq D_{Q_\tau}$,
\[
  Q_{X\setminus N}=Q_X\qquad\text{and}\qquad\hull(X\setminus N)=\hull(X).
\]
\end{corollary}

\begin{proof}
$Q_X\subseteq Q_{X\setminus N}$ trivially. Let $q\in Q_{X\setminus N}$ and $c\coloneqq
(1-\tau)/\tau$. If $x\in X\cap N$, then $P_{Q_\tau}(x)>0$, so by \eqref{eq:Ptau}
$\min_ix_i>-c\,P(x)$, and $q\cdot x\ge\min_ix_i>-c\,P(x)$ because $q\in\Delta$. For
a.e.\ $x\in X\setminus N$, $q\cdot x\ge0$ and $P(x)\ge0$, so $q\cdot x\ge-c\,P(x)$.
Lemma~\ref{lem:contraction} gives $q\in Q_X$.
\end{proof}

\begin{theorem}\label{thm:H}
Let $X\in\M_U$ be incoherent. Then there is $N\in\CC_U$ such that
\begin{enumerate}[label=(\roman*)]
  \item $m(N)-m(X)\gg0$, and
  \item $N\subseteq\hull(X\setminus N)$ and $\mu(N\setminus X)>0$.
\end{enumerate}
In particular $\mu\bigl((N\setminus X)\cap\hull(X\setminus N)\bigr)=\mu(N\setminus X)
>0$, and (H) holds.
\end{theorem}

\begin{proof}
\emph{Step 1: a strict improvement inside the hull.} Let $\M_G\coloneqq\{m(Y):Y
\subseteq G\text{ measurable}\}$. The vector measure $Y\mapsto\int_Yh\,d\mu$ on the
measurable subsets of $G$ is finite, by the first-moment condition in (M), and
non-atomic, since $\mu$ vanishes on hyperplanes and is therefore non-atomic. By Lyapunov's theorem
\cite{Lyapunov,Lindenstrauss}, $\M_G$ is compact and convex. Since $X\subseteq G$ a.e.,
$x\coloneqq m(X)\in\M_G$.

For $\lambda\in\Delta$ and measurable $Y\subseteq G$,
\[
  \lambda\cdot m(Y)=\int_Y\lambda\cdot h\,d\mu\le\int_{G\cap\{\lambda\cdot h>0\}}
  \lambda\cdot h\,d\mu ,
\]
with equality if and only if $Y=G\cap\{\lambda\cdot h>0\}$ a.e., because the hyperplane
$\{\lambda\cdot h=0\}$ is null. The maximiser is coherent: by Lemma~\ref{lem:D}(a),
\[
  G\cap\{\lambda\cdot h>0\}=D_{Q_X\cup\{\lambda\}}=D_{\conv(Q_X\cup\{\lambda\})},
\]
and $\conv(Q_X\cup\{\lambda\})$ is a credal set containing $K'$.

The function $f(\lambda,z)\coloneqq\lambda\cdot(z-x)$ is bilinear on the product of the
compact convex sets $\Delta$ and $\M_G$. By the minimax theorem \cite{Sion},
\[
  v\coloneqq\max_{z\in\M_G}\min_{\lambda\in\Delta}f(\lambda,z)
  =\min_{\lambda\in\Delta}\max_{z\in\M_G}f(\lambda,z),
\]
and both optima are attained, at $z^\ast$ and $\lambda^\ast$ say. Note that
$\min_{\lambda\in\Delta}f(\lambda,z)=\min_i(z_i-x_i)$.

Suppose $v\le0$. Then $\max_{z\in\M_G}\lambda^\ast\cdot z\le\lambda^\ast\cdot x$, so $X$
maximises $\lambda^\ast\cdot m(Y)$ over measurable $Y\subseteq G$, and $X=G\cap
\{\lambda^\ast\cdot h>0\}$ a.e. Then $X$ is coherent, a contradiction. Hence $v>0$.

Since $f(\lambda^\ast,z^\ast)\le\max_zf(\lambda^\ast,z)=v$ and $f(\lambda^\ast,z^\ast)
\ge\min_\lambda f(\lambda,z^\ast)=v$, the point $z^\ast$ maximises $\lambda^\ast\cdot z$
over $\M_G$. Writing $z^\ast=m(Y^\ast)$, uniqueness of the maximiser gives
$Y^\ast=N_0$ a.e., where
\[
  N_0\coloneqq G\cap\{\lambda^\ast\cdot h>0\}\in\CC_U,\qquad
  m(N_0)-m(X)=z^\ast-x\ge v\one\gg0 .
\]

\emph{Step 2: the gain is non-null; a uniform shrink.} Suppose $N_0\setminus X$ were
null. Then $m(N_0)-m(X)=-m(X\setminus N_0)$, and for $k\in K'$ we would have
$k\cdot(m(N_0)-m(X))=-\int_{X\setminus N_0}k\cdot h\,d\mu\le0$, since $k\cdot h>0$ on
$U$. This contradicts $m(N_0)-m(X)\gg0$ and $k\in\Delta$. So $\mu(N_0\setminus X)>0$.

For $0<\tau<1$ put $N_\tau\coloneqq N_0\cap D_{Q_\tau}$. By Lemma~\ref{lem:D}(a),
$N_\tau=D_{\conv(Q_X\cup\{\lambda^\ast\}\cup Q_\tau)}$, and this credal set contains $K'$;
so $N_\tau\in\CC_U$. Every $h\in N_0$ lies in $G$, so $P(h)>0$, and by \eqref{eq:Ptau}
$h\in D_{Q_\tau}$ for all small $\tau$. Hence $\mathbf 1_{N_\tau}\to\mathbf 1_{N_0}$
pointwise as $\tau\to0$, and dominated convergence, with dominating functions $|h|$ and
$1$, gives $m(N_\tau)\to m(N_0)$ and $\mu(N_0\setminus N_\tau)\to0$. Fix $\tau$ so small
that $m(N_\tau)-m(X)\gg0$ and $\mu(N_0\setminus N_\tau)<\mu(N_0\setminus X)$, and put
$N\coloneqq N_\tau$. Then (i) holds, and $\mu(N\setminus X)\ge\mu(N_0\setminus X)-
\mu(N_0\setminus N)>0$.

\emph{Step 3: the hull of what is removed.} Since $N\subseteq D_{Q_\tau}$,
Corollary~\ref{cor:shrink} gives $\hull(X\setminus N)=\hull(X)=G\supseteq N_0\supseteq
N$. This proves (ii), and $(N\setminus X)\cap\hull(X\setminus N)=N\setminus X$.
\end{proof}

\begin{remark}\label{rem:H}
\begin{enumerate}[label=(\alph*)]
  \item The proof uses no regularity of $X$, no pointedness of $\cl\hull(X)$, and no
  condition on the dimension of $K'$; nor does it use $K'\subseteq\relint\Delta$.
  \item The defender is explicit: $N=\hull(X)\cap\{\lambda^\ast\cdot h>0\}\cap
  D_{(1-\tau)Q_X+\tau\Delta}$, where $\lambda^\ast$ minimises the convex function
  $\lambda\mapsto\max_{z\in\M_G}\lambda\cdot z-\lambda\cdot m(X)$ over $\Delta$, and
  $\tau>0$ is small. Every density point $a$ of $N\setminus X$ lies in $N$: a density
  point of $N\setminus X$ is a density point of $N$; a point outside $\cl N$ is not one;
  and if $a\in\partial N$, a supporting half-space of the convex set $N$ at $a$ contains
  $N$, so $N$ fills at most half of every ball around $a$. Hence $a\in N\subseteq
  \hull(X\setminus N)$, and by Lemma~\ref{lem:hullformula} and Carath\'eodory's theorem,
  as in Corollary~\ref{cor:witness}, $a-\sum_{j=1}^r\lambda_jb_j\gg0$ for some $r\le n$,
  $b_j\in(X\setminus N)_d$ and $\lambda_j\ge0$.
  \item Step 1 applied with $G$ replaced by $U$ (which is coherent) shows that every
  $X\in\M_U$ that is not a.e.\ of the form $U\cap\{\pi\cdot h>0\}$, $\pi\in\Delta$, is
  uniformly dominated in Brier loss by a model of that form.
\end{enumerate}
\end{remark}

\section{Universal coverage}\label{sec:main}

\begin{theorem}\label{thm:main}
Let $K'\subseteq\Delta$ be nonempty, compact and convex, put $U\coloneqq
D_{K'}$, and consider the ranking $\CD(\CC_U)$ on $\M_U$.
\begin{enumerate}[label=(\alph*)]
  \item Each $\preceq_i$ is reflexive, ties a.e.-equal models, and satisfies (G), (IB),
  (TC) and (FCC). Each $\preceq_i$ is a sub-relation of the Brier ranking
  $\preceq^{\mathrm B}_i$ that is symmetric on Brier ties: $X\preceq_iY$ implies
  $m_i(X)\ge m_i(Y)$, and then $Y\preceq_iX$ if $m_i(X)=m_i(Y)$.
  \item The undominated models are exactly the coherent models on $U$:
  $\UU=\CC^\ast=\CC_U$. No coherent model has a weak dominator other than itself.
  \item Every incoherent $X\in\M_U$ is strictly dominated by a coherent undominated model
  $N$, which moreover uniformly dominates $X$ in Brier loss: $m(N)-m(X)\gg0$.
  \item (GC), (ID) and (W) hold, and (C) holds in the strong form: a model tied at every
  world with a coherent model equals it.
  \item Let $K$ be a credal set. $\CD(\CC_U)$ covers $K$ if and only if $K$ is
  full-dimensional and $K'\subseteq\relint K$.
\end{enumerate}
\end{theorem}

\begin{proof}
(a) is Theorem~\ref{thm:reduction}(a). (b) and (d): by Theorem~\ref{thm:H}, (H) holds,
so Corollary~\ref{cor:H} and Theorem~\ref{thm:reduction}(b),(e) apply with $\NN=\CC_U$.
(c): Theorem~\ref{thm:H} gives $N\in\CC_U$ with $m(N)-m(X)\gg0$ and $\mu((N\setminus X)
\cap\hull(X\setminus N))>0$. By Corollary~\ref{cor:H}, $r(N,X)$ attacks no coherent
model, so ($\star$) holds, and Theorem~\ref{thm:reduction}(c) shows that $N$ strictly
dominates $X$. $N$ is undominated by (b).

(e) Requirements (i), (ii), (iii) and (v) of Definition~\ref{def:cover} hold by (a) and
(d). By Lemma~\ref{lem:LF}(a), $D_K\in\M_U$ if and only if $K'\subseteq K$, and then
$D_K$ is coherent, hence undominated by (b). So $\CD(\CC_U)$ covers $K$ if and only if
$K'\subseteq K$ and (LF) holds at $D_K$. If $K$ is not full-dimensional, (LF) fails by
Lemma~\ref{lem:LF}(c), since $U$ is one-sided. If $K$ is full-dimensional and
$K'\subseteq K$, (LF) holds if and only if $K'\subseteq\relint K$, by
Lemma~\ref{lem:LF}(b).
\end{proof}

Since $\relint K\ne\emptyset$ exactly when $K$ is full-dimensional, the condition in (e)
can be written as $K'\subseteq\relint K$. It forces $K'\subseteq\relint\Delta$; if $K'$
meets the relative boundary of $\Delta$, no credal set is covered, although (a)--(d) still
hold.

\begin{definition}\label{def:universal}
Let $V$ be a domain. A credal set $K$ is \emph{coverable on $V$} if some ranking on $\M_V$
covers $K$. A ranking on $\M_V$ is a \emph{universal cover on $V$} if it covers every
credal set that is coverable on $V$.
\end{definition}

\begin{corollary}\label{cor:universal}
Let $K$ be a credal set. $\CD(\CC_U)$ covers $K$ if and only if $K$ is coverable on $U$.
Hence $\CD(\CC_U)$ is a universal cover on $U$.
\end{corollary}

\begin{proof}
If some ranking on $\M_U$ covers $K$, then $D_K\in\M_U$ and (LF) holds at $D_K$. So
$K'\subseteq K$ by Lemma~\ref{lem:LF}(a), $K$ is full-dimensional by
Lemma~\ref{lem:LF}(c), and $K'\subseteq\relint K$ by Lemma~\ref{lem:LF}(b). Now apply
Theorem~\ref{thm:main}(e). The converse is trivial.
\end{proof}

\begin{corollary}\label{cor:every}
Let $K$ be a full-dimensional credal set, and let $K'\subseteq\relint K$ be nonempty,
compact and convex; for instance $K'=\{k_0\}$ with $k_0\in\relint K$. Then
$\CD(\CC_{D_{K'}})$ covers $K$. In particular every full-dimensional credal set is
covered by some ranking on some domain. The same ranking covers every full-dimensional
credal set whose relative interior contains $K'$.
\end{corollary}

\begin{proof}
Theorem~\ref{thm:main}(e).
\end{proof}

\begin{proposition}\label{prop:lowdim}
Let $K$ be a credal set that is not full-dimensional, and let $V$ be a one-sided domain.
Then no ranking on $\M_V$ covers $K$.
\end{proposition}

\begin{proof}
(LF) fails at $D_K$ by Lemma~\ref{lem:LF}(c).
\end{proof}

The same conclusion holds on every domain $V$ for which $V\cap(-V)$ contains no non-null
set of gambles with a positive and a negative coordinate
\cite[Prop.~2.10]{KCM}. Lower-dimensional targets therefore require domains that are
not one-sided, and on such domains some of them are covered by other rankings; for
instance, on the full domain the Brier ranking covers every precise model $\{p\}$ with
$p\in\relint\Delta$ \cite[Cor.~5.9]{KCM}. Whether the analogue of
Theorem~\ref{thm:main} holds on domains that are not one-sided is open. The construction of
Section~\ref{sec:deletion} still makes sense there, since a coherent model contains no
$h\le0$; but $Q_B$ may then be empty, and Step 2 of the proof of Theorem~\ref{thm:H} uses
an element of $K'$.

\section{Scope and limits}\label{sec:scope}

\subsection{Universality is relative to a fixed domain}

Theorem~\ref{thm:main} provides one ranking for each domain $U=D_{K'}$. That ranking makes
every coherent model on $U$ undominated and covers exactly the full-dimensional credal
sets whose relative interiors contain $K'$. For any ranking on $\M_{D_{K'}}$, a cover of
a full-dimensional $K$ requires $K'\subseteq\relint K$ (Lemma~\ref{lem:LF}(a),(b)); so two
full-dimensional credal sets with disjoint relative interiors are not covered on a common
domain of this form. Lower-dimensional credal sets are covered on no one-sided domain
(Proposition~\ref{prop:lowdim}).

\subsection{Coherence enters the definition of the ranking}

Whether $\CD(\CC_U)$ compares two models at world $i$ depends on whether the record of the
comparison attacks a coherent model. The definition of the ranking therefore refers to
the class of coherent models. Theorem~\ref{thm:main} shows that the coherent models on $U$ can be
represented as the undominated models of a ranking satisfying (G), (IB), (TC) and (FCC), with every incoherent model strictly dominated. It does not derive coherence
from accuracy considerations that are specified independently of coherence. The
dominance part of the theorem does not follow from the definition: by
Corollary~\ref{cor:H}, (ID) for $\CD(\CC_U)$ is equivalent to (H), which is
Theorem~\ref{thm:H}. Whether a comparison is deleted at world $i$ also depends on the
moments of the record at the other worlds, through $\RR(\CC_U)$ and the choice of $i(r)$.

\subsection{Transitivity and disjoint additivity}

Consider the following conditions at world $i$.
\begin{itemize}
  \item[(Tr)] \emph{Transitivity.} If $X\preceq_iY$ and $Y\preceq_iZ$, then $X\preceq_iZ$.
  \item[(Add)] \emph{Disjoint additivity.} If $X\preceq_iY$ and $Y\preceq_iZ$, and
  $(X\triangle Y)\cap(Y\triangle Z)$ is null, then $X\preceq_iZ$.
\end{itemize}
(Tr) implies (Add). For records $r_1=(A_1,B_1)$ and $r_2=(A_2,B_2)$ whose supports
$A_1\cup B_1$ and $A_2\cup B_2$ are disjoint, put $r_1\sqcup r_2\coloneqq(A_1\cup A_2,
B_1\cup B_2)$. If $(X\triangle Y)\cap(Y\triangle Z)$ is null, then $r(X,Z)=r(X,Y)\sqcup
r(Y,Z)$: a gamble in $X\triangle Y$ is not in $Y\triangle Z$, so it lies in $Z$ exactly
when it lies in $Y$, and symmetrically. So, for a relation given by a set $\Acc_i$ of
records, (Add) holds whenever $\Acc_i$ is closed under $\sqcup$.

The ranking $\CD(\CC_U)$ satisfies neither (Tr) nor (Add) in general
(Example~\ref{ex:add}). This is a property of that ranking, not of universal coverage:
Proposition~\ref{prop:lex} gives a deletion ranking that satisfies (Add) and (FCC), has
the properties (b)--(e) of Theorem~\ref{thm:main}, and is a universal cover on $U$.

\begin{example}\label{ex:add}
Let $n=3$ and let $\mu$ have density $e^{-|h|_1}$, which satisfies (M). Put
\[
  v_1\coloneqq\bigl(\tfrac15,\tfrac25,\tfrac25\bigr),\quad
  v_2\coloneqq\bigl(\tfrac25,\tfrac15,\tfrac25\bigr),\quad
  v_3\coloneqq\bigl(\tfrac25,\tfrac25,\tfrac15\bigr),\quad
  u\coloneqq\tfrac13\one,
\]
$K\coloneqq\conv\{v_1,v_2,v_3\}$, $K'\coloneqq\tfrac12(K+u)$ and $U\coloneqq D_{K'}$. Let
$N\coloneqq D_Q$ with $Q\coloneqq\conv\{\tfrac{20}{23}v_k+\tfrac3{23}u:k=1,2,3\}$. Then
$K'\subseteq Q\subseteq K$, so $D_K\subseteq N$ and $N,D_K\in\CC_U$. Put
$D\coloneqq N\setminus D_K$, $w\coloneqq(76,-43,157)$,
\[
  S_1\coloneqq D\cap\bigl(\{v_2\cdot h<0\}\cup\{v_1\cdot h<0<v_2\cdot h,\ w\cdot h>0\}
  \bigr),\qquad S_2\coloneqq D\setminus S_1 .
\]
Exact computation of the cone moments in rational arithmetic gives
\begin{align*}
  m(D)&=\tfrac{68269}{2370816}\,\one,\\
  m(S_1)&=\Bigl(\tfrac{1691077}{59270400},\ \tfrac{122655769}{2963520000},\
  -\tfrac{58937269}{2963520000}\Bigr)\approx(0.02853,\ 0.04139,\ -0.01989),\\
  m(S_2)&=\Bigl(\tfrac{163}{617400},\ -\tfrac{37319519}{2963520000},\
  \tfrac{16030391}{329280000}\Bigr)\approx(0.00026,\ -0.01259,\ 0.04868),
\end{align*}
and $\mu(N\cap\{v_2\cdot h<0,\ h_1<0\})=\tfrac{1853}{70560}>0$. Consider the models
\[
  Y_0\coloneqq D_K,\qquad Y_1\coloneqq D_K\cup S_2,\qquad Y_2\coloneqq N=D_K\cup S_1\cup
  S_2 ,
\]
and any family $\NN\subseteq\CC_U$ with $D_K\in\NN$, for instance $\NN=\CC_U$.
\begin{itemize}
  \item $r(Y_2,Y_1)=(S_1,\emptyset)$ has $m_1>0$, so it lies in $\Br_1$. It is not in
  $\RR(\NN)$, since $m_3(S_1)<0$, and it is not the inverse of a member of $\RR(\NN)$,
  since $m_1(S_1)>0$. So $Y_2\preceq_1Y_1$.
  \item In the same way, using $m_1(S_2)>0$ and $m_2(S_2)<0$, $Y_1\preceq_1Y_0$.
  \item $r(Y_2,Y_0)=(D,\emptyset)$ has $m(D)\gg0$ and attacks $D_K\in\NN$, so it lies in
  $\RR(\NN)$. It is not in $G_1$, since $D\supseteq N\cap\{v_2\cdot h<0\}$ meets
  $\{h_1<0\}$ in positive measure. So $i((D,\emptyset))=1$, and not $Y_2\preceq_1Y_0$.
\end{itemize}
Thus (Tr) fails at world 1, and since $(Y_2\triangle Y_1)\cap(Y_1\triangle Y_0)=S_1\cap
S_2=\emptyset$, so does (Add). Also $r(Y_0,Y_2)=(\emptyset,D)\notin\Br_1$, so $Y_0$ and
$Y_2$ are incomparable at world 1, consistently with (FCC).
\end{example}

The following deletion ranking satisfies (Add). Let $\NN\subseteq\CC_U$ be nonempty, and let $\mathrm{NA}$ be the set of records that
attack no member of $\NN$; the null record attacks every model, so it is not in
$\mathrm{NA}$. For a record $r$ put $L_i(r)\coloneqq\bigl(m_i(r),\sum_jm_j(r)\bigr)\in
\R^2$, and order $\R^2$ lexicographically: $(a,b)\ge_{\mathrm{lex}}(a',b')$ if $a>a'$,
or $a=a'$ and $b\ge b'$; $>_{\mathrm{lex}}$ is its strict part. Define the ranking $\mathrm{LD}(\NN)$ on $\M_U$ by
$X\preceq_iY$ iff $r(X,Y)\in\Acc^{\mathrm{lex}}_i$, where
\[
  \Acc^{\mathrm{lex}}_i\coloneqq G_i\cup\{r\in\mathrm{NA}:L_i(r)>_{\mathrm{lex}}0\}
  \cup\{r\in\mathrm{NA}:r^{-1}\in\mathrm{NA},\ L_i(r)=0\}.
\]
World $i$ compares by its Brier score, breaking ties by the total Brier score, but it
accepts a record that attacks a member of $\NN$ only if the record is in $G_i$.

\begin{proposition}\label{prop:lex}
Let $\NN\subseteq\CC_U$ be nonempty, and consider $\mathrm{LD}(\NN)$.
\begin{enumerate}[label=(\alph*)]
  \item each $\preceq_i$ is reflexive, ties a.e.-equal models, and satisfies (G), (IB),
  (FCC), (TC) and (Add); and $X\preceq_iY$ implies $m_i(X)\ge m_i(Y)$;
  \item no member of $\NN$ has a weak dominator other than itself;
  \item for $X\in\M_U$ and $N\in\NN$, $N$ dominates $X$ if and only if ($\star$) holds,
  and then $N$ strictly dominates $X$;
  \item for $\NN=\CC_U$, $\mathrm{LD}(\CC_U)$ satisfies (ID) if and only if (H) holds.
  Hence, by Theorem~\ref{thm:H}, its undominated models are exactly the coherent models on
  $U$, every incoherent model is strictly dominated by one of them, (GC), (C) in the
  strong form, (ID) and (W) hold, and $\mathrm{LD}(\CC_U)$ covers a credal set $K$ if
  and only if $K$ is full-dimensional and $K'\subseteq\relint K$. So $\mathrm{LD}(\CC_U)$
  is a universal cover on $U$.
\end{enumerate}
\end{proposition}

\begin{proof}
\emph{A fact.} If $r_1\sqcup r_2$ attacks $Z$, then $r_1$ and $r_2$ attack $Z$, since
$B_k\subseteq B_1\cup B_2$ and $A_k\subseteq A_1\cup A_2$. Hence if $r_1\in\mathrm{NA}$
then $r_1\sqcup s\in\mathrm{NA}$ for every $s$ whose support is disjoint from that of
$r_1$.

(a) The null record lies in $G_i$, so Lemma~\ref{lem:record} gives reflexivity, ties of
a.e.-equal models and (IB). A non-null record in $G_i$ has $m_i>0$, as in the proof of
Theorem~\ref{thm:reduction}(a). Its inverse has $L_i<_{\mathrm{lex}}0$ and is not in
$G_i$ (it would have to be null), so it is not accepted; this gives (G). Put
$\Phi_i(X)\coloneqq(m_i(X),\sum_jm_j(X))$, so that $L_i(r(X,Y))=\Phi_i(X)-\Phi_i(Y)$.
Every accepted record has $L_i\ge_{\mathrm{lex}}0$. An accepted record with $L_i=0$ is
null or lies in the third set, which is closed under inverses; so its inverse is
accepted. Lemma~\ref{lem:deletion}, applied with $\varphi\coloneqq-\Phi_i$ and $T=\R^2$
ordered by $\ge_{\mathrm{lex}}$ (so that $X\preceq_iY$ implies $-\Phi_i(X)\le_{\mathrm{lex}}
-\Phi_i(Y)$), gives (FCC) and (TC). The implication $X\preceq_iY
\Rightarrow m_i(X)\ge m_i(Y)$ holds because $L_i(r)\ge_{\mathrm{lex}}0$ forces
$m_i(r)\ge0$.

\emph{(Add).} By the discussion above it suffices that $\Acc^{\mathrm{lex}}_i$ is closed
under $\sqcup$. Let $r_1,r_2$ be accepted with disjoint supports. If both lie in $G_i$, so
does $r_1\sqcup r_2$. Otherwise one of them, say $r_1$, is a non-null member of
$\mathrm{NA}$, and $r_1\sqcup r_2\in\mathrm{NA}$ by the fact. $L_i$ is additive, so
$L_i(r_1\sqcup r_2)=L_i(r_1)+L_i(r_2)\ge_{\mathrm{lex}}0$. If the sum is
$>_{\mathrm{lex}}0$, the union is accepted. If it is $0$, both terms are $0$, because
$\R^2$ with the lexicographic order is an ordered group. Then $r_1$ lies in the third
set, so $r_1^{-1}\in\mathrm{NA}$, and $(r_1\sqcup r_2)^{-1}=r_1^{-1}\sqcup r_2^{-1}\in
\mathrm{NA}$ by the fact. So the union lies in the third set.

(b) Let $Y$ weakly dominate $N\in\NN$. The record $r(Y,N)$ attacks $N$, so it is not in
$\mathrm{NA}$, and being accepted at every world it lies in $G_i$ for every $i$. Then
$Y\setminus N\subseteq\Rpp$ a.e.\ and $N\setminus Y\subseteq\{h:h_i<0\text{ for all
}i\}$ a.e. The latter set misses $U$, and $\Rpp\subseteq N$; so $Y=N$ a.e.

(c) Let $r\coloneqq r(N,X)$. Suppose ($\star$). Then $r\in\mathrm{NA}$, $m(r)\in\Rp$
and $\sum_jm_j(r)>0$, so $L_i(r)>_{\mathrm{lex}}0$ for every $i$, and $r$ is accepted at
every world. The record $r^{-1}$ is non-null, has $L_i<_{\mathrm{lex}}0$, and so is
neither in $G_i$ nor in the other two sets. So $N\prec_iX$ for every $i$. Conversely, let
$N$ dominate $X$; then $r$ is non-null and accepted at every world. If $r\notin
\mathrm{NA}$, then $r\in G_i$ for every $i$, so $r=(A,\emptyset)$ a.e.\ with $A\subseteq
\Rpp$ non-null; such a record attacks no coherent model, since every coherent model
contains $\Rpp$, contradicting $r\notin\mathrm{NA}$. So $r\in\mathrm{NA}$, and
$L_i(r)\ge_{\mathrm{lex}}0$ gives $m(N)\ge m(X)$. If $m(N)=m(X)$, then $L_i(r)=0$ for
every $i$, $r$ is not in $G_i$, so $r$ lies in the third set at every world; then
$r^{-1}$ is accepted at every world, and $N$ does not dominate $X$. So ($\star$) holds.

(d) By Corollary~\ref{cor:H}, for $\NN=\CC_U$ condition ($\star$) for $X$ and $N$ is the
condition on $N$ in (H). If (ID) holds, (c) gives (H). If (H) holds, (c) shows that every
incoherent model is strictly dominated by a member of $\CC_U$, and the remaining claims
follow from (b) and (c) exactly as in the proofs of Theorem~\ref{thm:reduction}(e) and
Theorem~\ref{thm:main}(e). (H) holds by Theorem~\ref{thm:H}. The universal-cover claim
then follows as in Corollary~\ref{cor:universal}.
\end{proof}

The equivalence in (d) concerns $\mathrm{LD}(\CC_U)$ only; it does not say that every
ranking satisfying (Add) and covering some $K$ yields (H). We do not claim that
$\mathrm{LD}(\CC_U)$ is transitive.

\subsection{Relation to the Brier ranking}

By Theorem~\ref{thm:main}(a), every comparison made by $\CD(\CC_U)$ at world $i$ is made
by the Brier ranking at world $i$, and every strict comparison is strict for the Brier
ranking; $\CD(\CC_U)$ is obtained from the Brier ranking by deleting comparisons. It is
not complete. In Example~\ref{ex:add}, the coherent model $N$ uniformly dominates the
coherent model $D_K$ in Brier loss, since $m(N)-m(D_K)=m(D)\gg0$, while both are
undominated in $\CD(\CC_U)$. More generally, by Remark~\ref{rem:H}(c), every coherent
model on $U$ that is not a.e.\ of the form $U\cap\{\pi\cdot h>0\}$ is uniformly dominated
in Brier loss by a model of that form, and is undominated in $\CD(\CC_U)$ by
Theorem~\ref{thm:main}(b).

\begin{thebibliography}{9}

\bibitem{KCM}
J.~Konek and C.~Campbell-Moore.
\newblock Covering coherent sets of desirable gambles.
\newblock Draft, 2026.

\bibitem{Lindenstrauss}
J.~Lindenstrauss.
\newblock A short proof of {L}iapounoff's convexity theorem.
\newblock \emph{Journal of Mathematics and Mechanics}, 15:971--972, 1966.

\bibitem{Lyapunov}
A.~Lyapunov.
\newblock Sur les fonctions-vecteurs compl\`etement additives.
\newblock \emph{Izvestiya Akademii Nauk SSSR, Seriya Matematicheskaya}, 4:465--478,
1940.

\bibitem{Rockafellar}
R.~T. Rockafellar.
\newblock \emph{Convex Analysis}.
\newblock Princeton University Press, 1970.

\bibitem{Sion}
M.~Sion.
\newblock On general minimax theorems.
\newblock \emph{Pacific Journal of Mathematics}, 8:171--176, 1958.

\end{thebibliography}

\end{document}
